% concatenated from Exact_Traces_JEE.tex
\RequirePackage{fix-cm}
\documentclass[smallextended,envcountsect,envcountsame]{svjour3}
\smartqed
\usepackage[T1]{fontenc}
\usepackage[utf8]{inputenc}
\usepackage{amsmath,amssymb,mathtools,mathrsfs,bm}
\usepackage{enumitem}
\usepackage{microtype}
\usepackage{hyperref}
\usepackage[nameinlink,capitalize]{cleveref}
\hypersetup{colorlinks=true,linkcolor=blue,citecolor=blue,urlcolor=blue}
\newcommand{\R}{\mathbb R}
\newcommand{\N}{\mathbb N}
\newcommand{\Kfun}{K}
\newcommand{\Tr}{\operatorname{Tr}}
\journalname{Journal of Evolution Equations}

\begin{document}
\title{Exact Traces for Vector-Valued Morrey Spaces}
\titlerunning{Exact Traces for Vector-Valued Morrey Spaces}
\author{Rishad Shahmurov \and Veli Shahmurov\thanks{Corresponding author}}
\authorrunning{R. Shahmurov and V. Shahmurov}
\institute{R. Shahmurov \at
Cellular Products Research and Development, 595 East Crossville Road, Roswell, GA 30075, USA \\
\email{rshahmurov@crimson.ua.edu}
\and
V. Shahmurov \at
Antalya Bilim University, Antalya, Turkey \\
\email{veli.sahmurov@antalya.edu.tr}}
\date{}
\maketitle

\begin{abstract}
Trace spaces determine exactly which initial values are compatible with an evolution class. We identify the exact trace generated by vector-valued Morrey control in time and show that it differs essentially from the classical $L^p$ theory. For a Banach couple $X_1\hookrightarrow X_0$, the trace of the natural Morrey evolution class is the weak real-interpolation space $(X_0,X_1)_{\theta,\infty}$, where $\theta=1-(1-\lambda)/p$. Every element of this space occurs as a trace through a bounded extension operator. The result is sharp in both interpolation parameters: in general the smoothness exponent cannot be increased and the fine index $\infty$ cannot be replaced by any finite index. Thus Morrey control changes the exact trace mechanism rather than merely strengthening an integrability estimate. At the limiting endpoint, bounded mean oscillation (BMO) control yields finite-index interpolation traces together with a logarithmic modulus of continuity. The result isolates the precise initial-data space naturally associated with local, scale-sensitive time regularity and provides a trace framework suited to evolution equations with Morrey-type maximal regularity.
\keywords{Morrey spaces \and trace spaces \and real interpolation \and evolution equations \and bounded mean oscillation}
\subclass{Primary 46E35, 47D06 \and Secondary 46B70, 46M35, 35K90}
\end{abstract}

\section{Introduction}
Trace spaces describe exactly which initial values are compatible with a given evolution class. In the classical $L^p$ setting, the theorem of Lions and Peetre identifies the trace of
\[
W^{1,p}(0,T;X_0)\cap L^p(0,T;X_1),\qquad X_1\hookrightarrow X_0,
\]
with $(X_0,X_1)_{1-1/p,p}$ and provides a bounded extension operator~\cite{LionsPeetre1964,Peetre1966}. This result became a basic ingredient of abstract parabolic theory and modern maximal regularity~\cite{Weis2001,KunstmannWeis2004}.

Morrey control changes the question. Instead of measuring the time variable only through one global $L^p$ norm, it controls the solution on every time interval with a scale-dependent normalization. It is therefore not clear a priori whether the classical trace survives, whether only the smoothness exponent changes, or whether the interpolation fine index changes as well.

We show that the change is structural. For $1<p<\infty$ and $0<\lambda<1$, set
\[
\theta=1-\frac{1-\lambda}{p}.
\]
For the natural Morrey class in which $u$ takes values in $X_1$ and $u'$ takes values in $X_0$, the trace at $t=0$ is exactly
\[
(X_0,X_1)_{\theta,\infty}.
\]
Every element of this space occurs as a trace, with a bounded extension back to the evolution class. Both parts of the answer are sharp: the exponent $\theta$ cannot be increased in general, and the weak fine index $\infty$ cannot be replaced by any finite $q$. In particular, passing from $L^p$ to Morrey control changes not only the trace exponent but also the fine interpolation scale.

Vector-valued trace theory in Sobolev, Besov, Triebel--Lizorkin, anisotropic and mixed-norm settings is extensive~\cite{Schneider2010,MeyriesVeraar2014,JohnsenSickel2018,AgrestiLindemulderVeraar2023,DenkRoodenburg2026}. Morrey interpolation and Besov--Morrey scales have also been studied~\cite{BurenkovEtAl2010,BurenkovGhorbanalizadehSawano2019,Peetre1969Morrey,Sawano2018}, and Morrey-type estimates occur in parabolic regularity~\cite{OgawaSuguro2023}. These results do not identify the exact trace space above, nor the forced weak fine index.

The endpoints reveal a second feature of the theory. The case $\lambda=1$ produces a subcritical family rather than the same exact identity, while $BMO$ control leads naturally to finite interpolation traces and, under a mixed $L^\infty$--$BMO$ hypothesis, to logarithmic continuity~\cite{JohnNirenberg1961}. We keep the paper focused on this trace mechanism and its sharpness; singular-integral lifting and maximal regularity are treated separately.
\section{Background and notation}\label{sec:background}
	
	\subsection{The \texorpdfstring{$K$}{K}-functional and real interpolation}
	
	Let $(X_0,X_1)$ be a compatible Banach couple, that is, $X_0$ and $X_1$ are continuously embedded into a common Hausdorff topological vector space. For $x\in X_0+X_1$ and $r>0$, the Peetre $K$-functional is defined by
	\[
	K(r,x;X_0,X_1)
	:=
	\inf_{x=x_0+x_1}
	\bigl(
	\|x_0\|_{X_0}+r\|x_1\|_{X_1}
	\bigr),
	\qquad x_0\in X_0,\ x_1\in X_1.
	\]
	
	We recall the basic properties used throughout the paper.
	
	\begin{proposition}\label{prop:K-basic}
		Let $(X_0,X_1)$ be a Banach couple, let $x\in X_0+X_1$, and let $r>0$. Then:
		
		\begin{enumerate}[label=\textup{(\roman*)}]
			\item $K(r,\cdot;X_0,X_1)$ is a seminorm on $X_0+X_1$, and an equivalent norm on the quotient by the natural kernel;
			
			\item the map
			\[
			r\mapsto K(r,x;X_0,X_1)
			\]
			is nondecreasing on $(0,\infty)$;
			
			\item the map
			\[
			r\mapsto r^{-1}K(r,x;X_0,X_1)
			\]
			is nonincreasing on $(0,\infty)$;
			
			\item for every decomposition $x=x_0+x_1$ with $x_0\in X_0$ and $x_1\in X_1$,
			\[
			K(r,x;X_0,X_1)\le \|x_0\|_{X_0}+r\|x_1\|_{X_1};
			\]
			
			\item if $x\in X_0\cap X_1$, then
			\[
			K(r,x;X_0,X_1)\le \min\{\|x\|_{X_0},\,r\|x\|_{X_1}\}.
			\]
		\end{enumerate}
	\end{proposition}
	
	\begin{proof}
		All assertions are standard consequences of the definition. The monotonicity in \textup{(ii)} follows because increasing $r$ increases the quantity
		\[
		\|x_0\|_{X_0}+r\|x_1\|_{X_1}
		\]
		for every fixed decomposition. Assertion \textup{(iii)} is obtained by dividing the defining expression by $r$ and observing that
		\[
		r^{-1}K(r,x;X_0,X_1)
		=
		\inf_{x=x_0+x_1}
		\bigl(
		r^{-1}\|x_0\|_{X_0}+\|x_1\|_{X_1}
		\bigr),
		\]
		which is nonincreasing in $r$. The remaining statements are immediate.
	\end{proof}
	
	For $0<\theta<1$ and $1\le q\le\infty$, the real interpolation space $(X_0,X_1)_{\theta,q}$ consists of all $x\in X_0+X_1$ such that
	\[
	\|x\|_{(X_0,X_1)_{\theta,q}}
	:=
	\begin{cases}
		\left(
		\displaystyle\int_0^\infty
		\bigl(r^{-\theta}K(r,x;X_0,X_1)\bigr)^q
		\,\frac{dr}{r}
		\right)^{1/q},
		& 1\le q<\infty,\\[2ex]
		\displaystyle\sup_{r>0} r^{-\theta}K(r,x;X_0,X_1),
		& q=\infty,
	\end{cases}
	\]
	is finite.
	
	\begin{remark}\label{rem:index-infinity}
		The case $q=\infty$ is the natural target in the present paper. Indeed, Morrey-in-time estimates are based on a supremum over time intervals and therefore yield uniform bounds of the form
		\[
		\sup_{r>0} r^{-\theta}K(r,u(t);X_0,X_1)\le C,
		\]
		rather than integral bounds in the interpolation variable $r$. This is exactly the norm of the weak real interpolation space $(X_0,X_1)_{\theta,\infty}$.
	\end{remark}
	
	\subsection{Interpolation scales: positive examples and Morrey-type caveats}
	
	Many classical function spaces form interpolation scales under the real method. This is the case, in particular, for Lebesgue spaces, Sobolev spaces, Besov spaces, and large classes of Triebel--Lizorkin spaces. Trace theorems on these scales are central in elliptic and parabolic analysis.
	
	By contrast, one must be more careful with global Morrey spaces. Standard Morrey spaces do not, in general, form a real interpolation scale in the same clean way as the classical scales above. However, important positive results are available in several related settings:
	\begin{enumerate}[label=\textup{(\roman*)}]
		\item Sobolev--Morrey couples admit interpolation descriptions in terms of Besov--Morrey or Nikol'skii--Besov--Morrey spaces;
		\item local Morrey-type spaces often have better interpolation behavior than global Morrey spaces;
		\item generalized Besov--Morrey and Triebel--Lizorkin--Morrey scales provide natural endpoint spaces for traces and embeddings.
	\end{enumerate}
	
	\begin{remark}\label{rem:morrey-scale-caveat}
		The abstract theorem proved below does not require Morrey spaces themselves to form an interpolation scale. We work with an arbitrary Banach couple $(X_0,X_1)$ and use only the general real-interpolation $K$-method. The Morrey structure enters solely through the time norm.
	\end{remark}
	
	\subsection{Time-Morrey spaces}
	
	\begin{definition}\label{def:time-morrey}
		Let $X$ be a Banach space, let $1<p<\infty$, let $0\le \lambda<1$, and let $T>0$. We define
		\[
		\|f\|_{\mathcal M^{p,\lambda}(0,T;X)}
		:=
		\sup_{I\subset(0,T)}
		|I|^{-\lambda/p}
		\left(\int_I \|f(t)\|_X^p\,dt\right)^{1/p},
		\]
		where the supremum is taken over all intervals $I\subset(0,T)$.
	\end{definition}
	
	\begin{definition}\label{def:holder-space}
		Let $X$ be a Banach space and let $0<\beta\le 1$. We write
		\[
		C^{0,\beta}([0,T];X)
		:=
		\Bigl\{u\in C([0,T];X):
		[u]_{C^{0,\beta}([0,T];X)}<\infty\Bigr\},
		\]
		where
		\[
		[u]_{C^{0,\beta}([0,T];X)}
		:=
		\sup_{0\le s<t\le T}
		\frac{\|u(t)-u(s)\|_X}{|t-s|^{\beta}}.
		\]
		The norm is
		\[
		\|u\|_{C^{0,\beta}([0,T];X)}
		:=
		\sup_{0\le t\le T}\|u(t)\|_X
		+
		[u]_{C^{0,\beta}([0,T];X)}.
		\]
		For $\beta=1$ this is the usual Lipschitz space. In particular, the statement
		\[
		u\in C^{0,\theta}([0,T];X_0)
		\]
		means that $u$ has a continuous $X_0$-valued representative and satisfies
		\[
		\|u(t)-u(s)\|_{X_0}\le C|t-s|^\theta
		\qquad(0\le s<t\le T).
		\]
	\end{definition}
	
	
	\begin{remark}\label{rem:lambda-limits}
		Let $X$ be a Banach space.
		
		\begin{enumerate}[label=\textup{(\roman*)}]
			\item If $\lambda=0$, then
			\[
			\mathcal M^{p,0}(0,T;X)=L^p(0,T;X)
			\]
			with equality of norms. In this regime,
			\[
			\theta=1-\frac1p.
			\]
			
			\item If $0<\lambda<1$, then the Morrey norm is stronger than the corresponding local $L^p$ control on short intervals and yields
			\[
			\int_I \|f(t)\|_{X}\,dt
			\le
			|I|^{1-\frac{1-\lambda}{p}}
			\|f\|_{\mathcal M^{p,\lambda}(0,T;X)}
			\]
			for every interval $I\subset(0,T)$.
			
			\item The trace exponent
			\[
			\theta=1-\frac{1-\lambda}{p}
			\]
			is strictly increasing in $\lambda$. As $\lambda\uparrow1$, one has $\theta\uparrow1$ and the trace space approaches the almost-strong endpoint
			\[
			(X_0,X_1)_{\theta,\infty},
			\qquad \theta\nearrow1.
			\]
			
			\item The endpoint $\lambda=1$ should not be described as a direct extension of Theorem~\ref{thm:morrey-trace-evolution}. By the Lebesgue differentiation theorem,
			\[
			\mathcal M^{p,1}(0,T;X)=L^\infty(0,T;X)
			\]
			with equivalent norms. Thus the endpoint evolution class becomes
			\[
			u\in W^{1,\infty}(0,T;X_0)\cap L^\infty(0,T;X_1).
			\]
			This gives Lipschitz continuity in $X_0$ and an essential $X_1$-bound, but it does not give strong continuity into $X_1$.
			
			\item The correct endpoint statement is therefore subcritical: for every $0<\eta<1$,
			\[
			W^{1,\infty}(0,T;X_0)\cap L^\infty(0,T;X_1)
			\hookrightarrow
			C^{0,1-\eta}([0,T];(X_0,X_1)_{\eta,\infty}).
			\]
			The endpoint $\eta=1$, which would correspond to $X_1$, generally fails without additional compactness, right-continuity, weak-continuity, or compatibility assumptions. This endpoint is treated separately in Subsection~\ref{subsec:endpoint-lambda-one}.
		\end{enumerate}
	\end{remark}

	


	\section{The Morrey trace theorem}\label{sec:trace}
	
	\begin{definition}\label{def:morrey-evolution-class}
		Let $(X_0,X_1)$ be a Banach couple with continuous embedding
		\[
		X_1\hookrightarrow X_0.
		\]
		Fix
		\[
		1<p<\infty,\qquad 0\le \lambda<1,\qquad T>0,
		\qquad
		\theta:=1-\frac{1-\lambda}{p}\in(0,1).
		\]
		We define
		\[
		\mathbb E^{p,\lambda}(0,T;X_0,X_1)
		:=
		\Bigl\{
		u:\;
		u\in \mathcal M^{p,\lambda}(0,T;X_1),\;
		\partial_t u\in \mathcal M^{p,\lambda}(0,T;X_0)
		\Bigr\},
		\]
		where $\partial_t u$ is understood in the sense of $X_0$-valued distributions.
	\end{definition}
	
	We now prove the Morrey trace theorem.

\begin{theorem}\label{thm:morrey-trace-evolution}
		Let $(X_0,X_1)$, $p$, $\lambda$, $T$, and $\theta$ be as in Definition~\ref{def:morrey-evolution-class}. Let
		\[
		u\in \mathbb E^{p,\lambda}(0,T;X_0,X_1).
		\]
		Then $u$ admits an $X_0$-valued absolutely continuous representative on $[0,T]$, still denoted by $u$, such that:
		
		\begin{enumerate}[label=\textup{(\roman*)}]
			\item for every $0\le s<t\le T$,
			\[
			u(t)-u(s)=\int_s^t \partial_t u(\tau)\,d\tau
			\qquad\text{in }X_0,
			\]
			and
			\[
			\|u(t)-u(s)\|_{X_0}
			\le
			|t-s|^\theta
			\|\partial_t u\|_{\mathcal M^{p,\lambda}(0,T;X_0)};
			\]
			
			\item for every $t\in[0,T]$,
			\[
			u(t)\in (X_0,X_1)_{\theta,\infty},
			\]
			and
			\[
			\sup_{t\in[0,T]}
			\|u(t)\|_{(X_0,X_1)_{\theta,\infty}}
			\le
			C\Bigl(
			\|u\|_{\mathcal M^{p,\lambda}(0,T;X_1)}
			+
			\|\partial_t u\|_{\mathcal M^{p,\lambda}(0,T;X_0)}
			\Bigr);
			\]
			
			\item in particular,
			\[
			u\in C^{0,\theta}([0,T];X_0)\cap L^\infty\bigl(0,T;(X_0,X_1)_{\theta,\infty}\bigr),
			\]
			and the endpoint traces
			\[
			u(0),\qquad u(T)
			\]
			are well-defined elements of $(X_0,X_1)_{\theta,\infty}$.
		\end{enumerate}
	\end{theorem}
	
	\begin{proof}
		We divide the proof into four steps.
		
		\medskip
		\noindent
		\textit{Step 1: Morrey control implies local $L^1$ control.}
		Let $X$ be a Banach space and let $f\in \mathcal M^{p,\lambda}(0,T;X)$. For any interval $I\subset(0,T)$, Hölder's inequality gives
		\[
		\int_I \|f(t)\|_X\,dt
		\le
		|I|^{1-\frac1p}
		\left(\int_I \|f(t)\|_X^p\,dt\right)^{1/p}.
		\]
		By Definition~\ref{def:time-morrey},
		\[
		\left(\int_I \|f(t)\|_X^p\,dt\right)^{1/p}
		\le
		|I|^{\lambda/p}\|f\|_{\mathcal M^{p,\lambda}(0,T;X)}.
		\]
		Hence
		\[
		\int_I \|f(t)\|_X\,dt
		\le
		|I|^{1-\frac{1-\lambda}{p}}
		\|f\|_{\mathcal M^{p,\lambda}(0,T;X)}
		=
		|I|^\theta
		\|f\|_{\mathcal M^{p,\lambda}(0,T;X)}.
		\]
		Therefore
		\[
		\mathcal M^{p,\lambda}(0,T;X)\hookrightarrow L^1(0,T;X)
		\]
		continuously on bounded intervals.
		
		Applying this with $X=X_0$ and using $X_1\hookrightarrow X_0$, we get
		\[
		u\in L^1(0,T;X_0),
		\qquad
		\partial_t u\in L^1(0,T;X_0).
		\]
		Thus
		\[
		u\in W^{1,1}(0,T;X_0),
		\]
		hence $u$ admits an absolutely continuous $X_0$-valued representative on $[0,T]$.
		
		\medskip
		\noindent
		\textit{Step 2: Hölder continuity in $X_0$.}
		Since $u\in W^{1,1}(0,T;X_0)$, for every $0\le s<t\le T$,
		\[
		u(t)-u(s)=\int_s^t \partial_t u(\tau)\,d\tau
		\qquad\text{in }X_0.
		\]
		Therefore
		\[
		\|u(t)-u(s)\|_{X_0}
		\le
		\int_s^t \|\partial_t u(\tau)\|_{X_0}\,d\tau
		\le
		|t-s|^\theta
		\|\partial_t u\|_{\mathcal M^{p,\lambda}(0,T;X_0)}.
		\]
		This proves \textup{(i)}.
		
		\medskip
		\noindent
		\textit{Step 3: Pointwise $K$-functional estimate.}
		Fix $t\in[0,T]$. For admissible $h>0$, define
		\[
		b_h(t):=\frac1h\int_t^{t+h}u(\tau)\,d\tau\in X_1,
		\qquad
		a_h(t):=u(t)-b_h(t)\in X_0,
		\]
		or the analogous backward average near $T$. Then
		\[
		u(t)=a_h(t)+b_h(t).
		\]
		
		Since
		\[
		u(\tau)-u(t)=\int_t^\tau \partial_t u(s)\,ds,
		\]
		we obtain
		\[
		\|a_h(t)\|_{X_0}
		\le
		h^\theta
		\|\partial_t u\|_{\mathcal M^{p,\lambda}(0,T;X_0)}.
		\]
		Also,
		\[
		\|b_h(t)\|_{X_1}
		\le
		h^{-(1-\theta)}
		\|u\|_{\mathcal M^{p,\lambda}(0,T;X_1)}.
		\]
		
		Thus, by Proposition~\ref{prop:K-basic},
		\[
		K(r,u(t);X_0,X_1)
		\le
		h^\theta A + r h^{-(1-\theta)} B,
		\]
		where
		\[
		A:=\|\partial_t u\|_{\mathcal M^{p,\lambda}(0,T;X_0)},
		\qquad
		B:=\|u\|_{\mathcal M^{p,\lambda}(0,T;X_1)}.
		\]
		We split the interpolation scales into two ranges.  If $0<r\le T/2$, then for every $t\in[0,T]$ at least one of the intervals $[t,t+r]$ and $[t-r,t]$ is contained in $[0,T]$.  Using the corresponding forward or backward average and choosing $h=r$ gives
		\[
		r^{-\theta}K(r,u(t);X_0,X_1)
		\le C(A+B),
		\qquad 0<r\le T/2,
		\]
		uniformly in $t$.

		For the remaining scales we use the $X_0$ part of the $K$-functional.  The Morrey bound on the full interval gives
		\[
		\left(\int_0^T\|u(s)\|_{X_1}^p\,ds\right)^{1/p}
		\le T^{\lambda/p}B.
		\]
		Hence there is a point $\tau\in(0,T)$ at which the $X_1$-valued representative is defined and
		\[
		\|u(\tau)\|_{X_1}
		\le T^{-(1-\lambda)/p}B
		=T^{\theta-1}B.
		\]
		Writing $c_{10}$ for the norm of the embedding $X_1\hookrightarrow X_0$, Step~2 yields
		\[
		\sup_{t\in[0,T]}\|u(t)\|_{X_0}
		\le c_{10}T^{\theta-1}B+T^\theta A.
		\]
		Therefore, if $r>T/2$,
		\[
		r^{-\theta}K(r,u(t);X_0,X_1)
		\le r^{-\theta}\|u(t)\|_{X_0}
		\le C_{T,c_{10}}(A+B).
		\]
		Combining the two scale ranges, we obtain
		\[
		\sup_{r>0}r^{-\theta}K(r,u(t);X_0,X_1)
		\le C_{p,\lambda,T,c_{10}}(A+B),
		\]
		uniformly in $t$.
		
		By the definition of the real interpolation norm,
		\[
		\|u(t)\|_{(X_0,X_1)_{\theta,\infty}}
		\le
		C(A+B).
		\]
		This proves \textup{(ii)}.
		
		\medskip
		\noindent
		\textit{Step 4: Endpoint traces.}
		Since $u\in C([0,T];X_0)$ by Step~2, the values $u(0)$ and $u(T)$ are well defined in $X_0$. Step~3 shows they belong in fact to $(X_0,X_1)_{\theta,\infty}$ and satisfy the same estimate. This proves \textup{(iii)}.
	\end{proof}
	
	\subsection{Exactness and sharpness of the Morrey trace space}\label{subsec:morrey-trace-sharpness}

	The preceding theorem gives the embedding direction of the Morrey trace theorem.  In the genuinely Morrey regime the target space is not only a convenient weak endpoint; it is the exact trace range.  The case \(\lambda=0\) is exceptional: then \(\mathcal M^{p,0}=L^p\), and the classical trace theorem identifies the sharper trace space \((X_0,X_1)_{1-1/p,p}\).  The theorem below therefore concerns \(0<\lambda<1\).

	We next identify the full trace range.

\begin{theorem}\label{thm:exact-morrey-trace-space}
	Let \((X_0,X_1)\) be a Banach couple with continuous embedding \(X_1\hookrightarrow X_0\).  Let
	\[
	1<p<\infty,\qquad 0<\lambda<1,\qquad T>0,
	\qquad
	\theta=1-\frac{1-\lambda}{p}.
	\]
	Then the trace at \(t=0\) of the Morrey evolution class is exactly
	\[
	\operatorname{Tr}_{0}\mathbb E^{p,\lambda}(0,T;X_0,X_1)
	=
	(X_0,X_1)_{\theta,\infty}.
	\]
	More precisely, the trace map
	\[
	\operatorname{Tr}_{0}:u\mapsto u(0)
	\]
	is bounded from \(\mathbb E^{p,\lambda}(0,T;X_0,X_1)\), equipped with its natural norm, into \((X_0,X_1)_{\theta,\infty}\).  Conversely, for every \(x\in (X_0,X_1)_{\theta,\infty}\) there exists
	\[
	u\in \mathbb E^{p,\lambda}(0,T;X_0,X_1)
	\]
	such that \(u(0)=x\) in \(X_0\), and
	\[
	\|u\|_{\mathcal M^{p,\lambda}(0,T;X_1)}
	+
	\|u'\|_{\mathcal M^{p,\lambda}(0,T;X_0)}
	\le
	C\|x\|_{(X_0,X_1)_{\theta,\infty}}.
	\]
	The constant \(C\) depends only on \(p,\lambda,T\), and the embedding constant of \(X_1\hookrightarrow X_0\).
	\end{theorem}

	\begin{proof}
	The boundedness of the trace map is exactly the trace estimate in Theorem~\ref{thm:morrey-trace-evolution}.  We prove the converse by an explicit dyadic extension.

	Let
	\[
	M:=\|x\|_{(X_0,X_1)_{\theta,\infty}}.
	\]
	For \(n=0,1,2,\ldots\), put
	\[
	r_n:=T2^{-n}.
	\]
	By the definition of the \((X_0,X_1)_{\theta,\infty}\)-norm, for each \(n\) we may choose a decomposition
	\[
	x=a_n+b_n,
	\qquad a_n\in X_0,
	\qquad b_n\in X_1,
	\]
	such that
	\[
	\|a_n\|_{X_0}+r_n\|b_n\|_{X_1}
	\le
	2K(r_n,x;X_0,X_1)
	\le
	2Mr_n^\theta.
	\]
	Therefore
	\[
	\|x-b_n\|_{X_0}=\|a_n\|_{X_0}\le 2Mr_n^\theta,
	\qquad
	\|b_n\|_{X_1}\le 2Mr_n^{\theta-1}.
	\]
	In particular \(b_n\to x\) in \(X_0\).

	Define \(u\) on each dyadic interval \((r_{n+1},r_n]\) by linear interpolation between \(b_{n+1}\) and \(b_n\):
	\[
	u(t):=
	\frac{t-r_{n+1}}{r_n-r_{n+1}}b_n
	+
	\frac{r_n-t}{r_n-r_{n+1}}b_{n+1},
	\qquad r_{n+1}<t\le r_n.
	\]
	Set \(u(0):=x\).  Since \(b_n\to x\) in \(X_0\), the resulting function is continuous at \(0\) as an \(X_0\)-valued function.  For \(t\in(r_{n+1},r_n]\), the preceding estimates give
	\[
	\|u(t)\|_{X_1}
	\le
	\max\{\|b_n\|_{X_1},\|b_{n+1}\|_{X_1}\}
	\le
	C M r_n^{\theta-1}.
	\]
	Since \(t\simeq r_n\) on \((r_{n+1},r_n]\), and \(\theta-1=-(1-\lambda)/p\), we have
	\[
	\|u(t)\|_{X_1}
	\le
	C M t^{-\frac{1-\lambda}{p}}.
	\]
	Similarly, on \((r_{n+1},r_n)\),
	\[
	u'(t)=\frac{b_n-b_{n+1}}{r_n-r_{n+1}}.
	\]
	Using \(b_n-x=-a_n\) and \(b_{n+1}-x=-a_{n+1}\),
	\[
	\|b_n-b_{n+1}\|_{X_0}
	\le
	\|a_n\|_{X_0}+\|a_{n+1}\|_{X_0}
	\le
	C M r_n^\theta.
	\]
	Since \(r_n-r_{n+1}=r_{n+1}\simeq r_n\), it follows that
	\[
	\|u'(t)\|_{X_0}
	\le
	C M r_n^{\theta-1}
	\le
	C M t^{-\frac{1-\lambda}{p}}.
	\]

	It remains only to observe that the scalar function
	\[
	g(t)=t^{-\frac{1-\lambda}{p}},\qquad 0<t<T,
	\]
	belongs to \(\mathcal M^{p,\lambda}(0,T)\) when \(0<\lambda<1\).  Indeed, if \(I=(a,b)\subset(0,T)\) and \(h=b-a\), then
	\[
	\int_I g(t)^p\,dt=\int_a^b t^{-(1-\lambda)}\,dt.
	\]
	If \(a\le h\), then this is bounded by \(C h^\lambda\).  If \(a>h\), then
	\[
	\int_a^b t^{-(1-\lambda)}\,dt
	\le
	h a^{-(1-\lambda)}
	\le
	h h^{-(1-\lambda)}=h^\lambda.
	\]
	Thus
	\[
	|I|^{-\lambda/p}
	\left(\int_I g(t)^p\,dt\right)^{1/p}
	\le C,
	\]
	uniformly in \(I\).  Hence \(u\in\mathcal M^{p,\lambda}(0,T;X_1)\) and \(u'\in\mathcal M^{p,\lambda}(0,T;X_0)\), with the asserted estimate.
	\end{proof}

	\begin{proposition}\label{prop:sharp-morrey-holder-exponent}
	The exponent
	\[
	\theta=1-\frac{1-\lambda}{p}
	\]
	in the \(X_0\)-valued H\"older conclusion of Theorem~\ref{thm:morrey-trace-evolution} cannot be improved.
	\end{proposition}

	\begin{proof}
	For \(0<\lambda<1\), take \(X_0=X_1=\mathbb R\) and
	\[
	f(t)=t^{-\frac{1-\lambda}{p}},
	\qquad
	u(t)=\int_0^t f(s)\,ds=C_{p,\lambda}t^\theta.
	\]
	As shown in the proof of Theorem~\ref{thm:exact-morrey-trace-space}, \(f\in\mathcal M^{p,\lambda}(0,1)\).  Hence \(u'\in\mathcal M^{p,\lambda}(0,1)\), while \(u\) is not \(\beta\)-H\"older at \(0\) for any \(\beta>\theta\).

	For \(\lambda=0\), the borderline power \(t^{-1/p}\) is not in \(L^p(0,1)\).  A logarithmic correction gives the same conclusion.  Let
	\[
	f(t)=t^{-1/p}\bigl(\log(e/t)\bigr)^{-1},
	\qquad 0<t<1.
	\]
	Then \(f\in L^p(0,1)=\mathcal M^{p,0}(0,1)\), because
	\[
	\int_0^1 t^{-1}\bigl(\log(e/t)\bigr)^{-p}\,dt<\infty
	\]
	for \(p>1\).  Its primitive satisfies, up to harmless constants,
	\[
	u(t)\simeq t^{1-1/p}\bigl(\log(e/t)\bigr)^{-1}
	\]
	as \(t\downarrow0\).  Therefore \(u\) is not \(\beta\)-H\"older at \(0\) for any \(\beta>1-1/p\).  This proves sharpness also in the classical \(\lambda=0\) endpoint of the Morrey scale.
	\end{proof}

	\begin{proposition}\label{prop:q-infty-forced}
	Let \(0<\lambda<1\), \(1<p<\infty\), and \(\theta=1-(1-\lambda)/p\).  In general, the trace target in Theorem~\ref{thm:morrey-trace-evolution} cannot be replaced by \((X_0,X_1)_{\theta,q}\) for any finite \(q\).
	\end{proposition}

	\begin{proof}
	By Theorem~\ref{thm:exact-morrey-trace-space}, it suffices to exhibit a Banach couple and a vector which belongs to \((X_0,X_1)_{\theta,\infty}\) but not to \((X_0,X_1)_{\theta,q}\) for any finite \(q\).

	Let
	\[
	X_0=c_0,
	\qquad
	X_1=\Bigl\{x=(x_n)_{n\ge1}:\sup_{n\ge1}2^n|x_n|<\infty\Bigr\},
	\]
	with the natural norms
	\[
	\|x\|_{X_0}=\sup_n |x_n|,
	\qquad
	\|x\|_{X_1}=\sup_n 2^n|x_n|.
	\]
	Then \(X_1\hookrightarrow X_0\).  For this diagonal couple, the scalar coordinate calculation gives
	\[
	K(r,x;X_0,X_1)
	\simeq
	\sup_{n\ge1}\min\{1,r2^n\}|x_n|.
	\]
	Now take
	\[
	x_n=2^{-\theta n}.
	\]
	Then \(x\in c_0\), and
	\[
	\sup_{r>0}r^{-\theta}K(r,x;X_0,X_1)<\infty.
	\]
	Indeed, if \(2^{-(N+1)}<r\le2^{-N}\), then
	\[
	\sup_n \min\{1,r2^n\}2^{-\theta n}
	\lesssim r^\theta.
	\]
	Hence \(x\in (X_0,X_1)_{\theta,\infty}\).

	On the other hand, for the same \(r\)-range one also has the lower estimate
	\[
	K(r,x;X_0,X_1)
	\gtrsim r^\theta.
	\]
	Consequently
	\[
	\int_0^1\bigl(r^{-\theta}K(r,x;X_0,X_1)\bigr)^q\,\frac{dr}{r}
	=\infty
	\]
	for every finite \(q\).  Thus \(x\notin (X_0,X_1)_{\theta,q}\) for \(q<\infty\).  Since \(x\) is nevertheless the trace of a function in \(\mathbb E^{p,\lambda}\) by Theorem~\ref{thm:exact-morrey-trace-space}, no finite interpolation index can replace \(\infty\) in general.
	\end{proof}

	\begin{remark}\label{rem:no-continuity-trace-space}
	The exact trace theorem identifies the range of the pointwise trace map.  It should not be read as a strong-continuity statement with values in \((X_0,X_1)_{\theta,\infty}\).  In general, Morrey control gives an \(X_0\)-continuous representative and uniformly bounded \((X_0,X_1)_{\theta,\infty}\)-traces.  Strong continuity in the weak interpolation norm requires additional hypotheses, such as compactness, stronger summability, or an independent continuity assumption in the interpolation space.
	\end{remark}

	

	\subsection{The endpoint \texorpdfstring{$\lambda=1$}{lambda=1} and BMO trace consequences}\label{subsec:endpoint-lambda-one}

	The endpoint $\lambda=1$ is not an endpoint of Theorem~\ref{thm:morrey-trace-evolution} in the literal sense $\theta=1$. It is a different theorem. With the normalization of Definition~\ref{def:time-morrey}, the endpoint Morrey space is $L^\infty$; hence the natural endpoint evolution class is
	\[
	\mathbb E^\infty(0,T;X_0,X_1)
	:=
	W^{1,\infty}(0,T;X_0)\cap L^\infty(0,T;X_1).
	\]
	The conclusion is continuity into every subcritical interpolation space, but not generally into $X_1$.

	\begin{lemma}\label{lem:morrey-endpoint-linfty}
	Let $X$ be a Banach space and $1<p<\infty$. If
	\[
	\|g\|_{\mathcal M^{p,1}(0,T;X)}
	:=
	\sup_{I\subset(0,T)} |I|^{-1/p}
	\left(\int_I\|g(t)\|_X^p\,dt\right)^{1/p},
	\]
	then
	\[
	\mathcal M^{p,1}(0,T;X)=L^\infty(0,T;X)
	\]
	with equivalent norms.
	\end{lemma}

	\begin{proof}
	The inclusion $L^\infty\subset\mathcal M^{p,1}$ is immediate. Conversely, the defining estimate gives uniform bounds for the $L^p$ averages of $\|g\|_X$ over all intervals. By the Lebesgue differentiation theorem applied to the scalar function $\|g(\cdot)\|_X^p$, these averages converge almost everywhere to $\|g(t)\|_X^p$, giving the $L^\infty$ bound.
	\end{proof}

	\begin{theorem}\label{thm:endpoint-subcritical-trace}
	Let $(X_0,X_1)$ be a Banach couple with $X_1\hookrightarrow X_0$. Assume
	\[
	u\in W^{1,\infty}(0,T;X_0)\cap L^\infty(0,T;X_1).
	\]
	Then, for every $0<\eta<1$,
	\[
	u\in C^{0,1-\eta}([0,T];(X_0,X_1)_{\eta,\infty}),
	\]
	and
	\begin{align*}
	&[u]_{C^{0,1-\eta}([0,T];(X_0,X_1)_{\eta,\infty})}
	\le
	C_\eta
	\|u'\|_{L^\infty(0,T;X_0)}^{1-\eta}
	\|u\|_{L^\infty(0,T;X_1)}^{\eta}.
	\end{align*}
	Moreover,
	\[
	\sup_{0\le t\le T}\|u(t)\|_{(X_0,X_1)_{\eta,\infty}}
	\le
	C_\eta
	\|u\|_{L^\infty(0,T;X_0)}^{1-\eta}
	\|u\|_{L^\infty(0,T;X_1)}^\eta.
	\]
	\end{theorem}

	\begin{proof}
	We first record the elementary interpolation inequality
	\[
	\|x\|_{(X_0,X_1)_{\eta,\infty}}
	\le
	C_\eta \|x\|_{X_0}^{1-\eta}\|x\|_{X_1}^{\eta},
	\qquad x\in X_0\cap X_1.
	\]
	Indeed,
	\[
	K(r,x;X_0,X_1)\le \min\{\|x\|_{X_0},r\|x\|_{X_1}\},
	\]
	and taking the supremum of $r^{-\eta}K(r,x)$ over $r>0$ gives the estimate.

	For almost every $s,t$ we may apply this to $x=u(t)-u(s)$. Since
	\[
	\|u(t)-u(s)\|_{X_0}
	\le |t-s|\|u'\|_{L^\infty(0,T;X_0)}
	\]
	and
	\[
	\|u(t)-u(s)\|_{X_1}
	\le 2\|u\|_{L^\infty(0,T;X_1)},
	\]
	we obtain
	\[
	\|u(t)-u(s)\|_{(X_0,X_1)_{\eta,\infty}}
	\le
	C_\eta |t-s|^{1-\eta}
	\|u'\|_{L^\infty(X_0)}^{1-\eta}
	\|u\|_{L^\infty(X_1)}^{\eta}.
	\]
	The estimate extends to every pair $s,t\in[0,T]$ by taking the $X_0$-continuous representative and approximating by Lebesgue points of the $X_1$-valued representative. The same interpolation inequality applied to $x=u(t)$ gives the uniform trace bound. This proves the theorem.
	\end{proof}

	\begin{remark}\label{rem:endpoint-no-x1}
	Theorem~\ref{thm:endpoint-subcritical-trace} reaches every subcritical space $(X_0,X_1)_{\eta,\infty}$, $0<\eta<1$. It does not generally reach $X_1$. The available endpoint information is only
	\[
	u\in L^\infty(0,T;X_1),
	\qquad
	u\in C^{0,1}([0,T];X_0),
	\]
	and these two facts do not imply $u\in C([0,T];X_1)$. Hence $\lambda=1$ is included through a different endpoint theorem rather than by inserting $\theta=1$ into the Morrey trace theorem.
	\end{remark}


	\subsubsection{BMO endpoint trace consequences}\label{subsubsec:bmo-trace}

	The endpoint of Calder\'on--Zygmund theory is naturally expressed in terms of \(BMO\).  The following trace consequences are therefore useful when maximal-regularity estimates give \(Au\) and \(u'\) in \(BMO\), rather than in \(L^\infty\).  We use the finite-interval norm
	\[
	\|g\|_{BMO_*(0,T;X)}:=\|g_{(0,T)}\|_X+[g]_{BMO(0,T;X)},
	\]
	where
	\[
	g_I:=\frac1{|I|}\int_I g(t)\,dt,
	\qquad
	[g]_{BMO(0,T;X)}:=
	\sup_{I\subset(0,T)}\frac1{|I|}\int_I\|g(t)-g_I\|_X\,dt.
	\]
	The average term is included because \(BMO\) itself is only a seminorm modulo constants.

	\begin{lemma}\label{lem:bmo-to-lq}
	Let \(X\) be a Banach space and \(1<q<\infty\).  Then
	\[
	BMO_*(0,T;X)\hookrightarrow L^q(0,T;X),
	\]
	and
	\[
	\|g\|_{L^q(0,T;X)}
	\le
	C_q T^{1/q}\|g\|_{BMO_*(0,T;X)}.
	\]
	\end{lemma}

	\begin{proof}
	Let \(J=(0,T)\) and write \(g_J\) for the average of \(g\) on \(J\).  Set
	\[
	h(t):=\|g(t)-g_J\|_X.
	\]
	We first show that \(h\) has scalar \(BMO\) seminorm controlled by the vector-valued \(BMO\) seminorm of \(g\).  For every interval \(I\subset J\), the scalar average \(h_I\) satisfies
	\[
	\frac1{|I|}\int_I |h(t)-h_I|\,dt
	\le
	\frac2{|I|}\int_I |h(t)-\|g_I-g_J\|_X|\,dt.
	\]
	By the reverse triangle inequality,
	\[
	\bigl|\|g(t)-g_J\|_X-\|g_I-g_J\|_X\bigr|
	\le
	\|g(t)-g_I\|_X.
	\]
	Therefore
	\[
	[h]_{BMO(0,T)}\le 2[g]_{BMO(0,T;X)}.
	\]
	Also
	\[
	h_J=\frac1T\int_0^T\|g(t)-g_J\|_X\,dt
	\le [g]_{BMO(0,T;X)}.
	\]
	The scalar John--Nirenberg inequality~\cite{JohnNirenberg1961} on the finite interval \(J\) gives
	\[
	\|h\|_{L^q(0,T)}
	\le
	C_q T^{1/q}\bigl(h_J+[h]_{BMO(0,T)}\bigr)
	\le
	C_q T^{1/q}[g]_{BMO(0,T;X)}.
	\]
	Finally,
	\[
	\|g\|_{L^q(0,T;X)}
	\le
	\|g-g_J\|_{L^q(0,T;X)}+T^{1/q}\|g_J\|_X
	\le
	C_qT^{1/q}\|g\|_{BMO_*(0,T;X)}.
	\]
	\end{proof}

	\begin{theorem}\label{thm:bmo-lions-peetre-trace}
	Let \((X_0,X_1)\) be a Banach couple with \(X_1\hookrightarrow X_0\).  Assume
	\[
	u\in BMO_*(0,T;X_1),
	\qquad
	u'\in BMO_*(0,T;X_0),
	\]
	where \(u'\) is the \(X_0\)-valued distributional derivative.  Then for every \(1<q<\infty\),
	\[
	u\in C\bigl([0,T];(X_0,X_1)_{1-1/q,q}\bigr),
	\]
	and
	\[
	\|u\|_{C([0,T];(X_0,X_1)_{1-1/q,q})}
	\le
	C_q\Bigl(
	\|u\|_{BMO_*(0,T;X_1)}+
	T\|u'\|_{BMO_*(0,T;X_0)}
	\Bigr),
	\]
	where $C_q$ is independent of $T$.
	\end{theorem}

	\begin{proof}
	By Lemma~\ref{lem:bmo-to-lq},
	\[
	u\in L^q(0,T;X_1),
	\qquad
	u'\in L^q(0,T;X_0).
	\]
	Hence
	\[
	u\in W^{1,q}(0,T;X_0)\cap L^q(0,T;X_1).
	\]
	Rescaling the classical trace theorem from $(0,1)$ to $(0,T)$ gives
	\[
	\|u\|_{C([0,T];(X_0,X_1)_{1-1/q,q})}
	\le C_q\Bigl(
	T^{-1/q}\|u\|_{L^q(0,T;X_1)}
	+T^{1-1/q}\|u'\|_{L^q(0,T;X_0)}
	\Bigr).
	\]
	Applying Lemma~\ref{lem:bmo-to-lq} to the two terms gives
	\[
	T^{-1/q}\|u\|_{L^q(X_1)}\le C_q\|u\|_{BMO_*(X_1)},
	\qquad
	T^{1-1/q}\|u'\|_{L^q(X_0)}\le C_qT\|u'\|_{BMO_*(X_0)},
	\]
	which is the asserted estimate.
	\end{proof}

	\begin{theorem}\label{thm:bmo-log-trace}
	Let \((X_0,X_1)\) be a Banach couple with \(X_1\hookrightarrow X_0\).  Assume
	\[
	u\in L^\infty(0,T;X_1),
	\qquad
	u'\in BMO_*(0,T;X_0).
	\]
	Then \(u\) has an \(X_0\)-valued continuous representative and, for every \(0<\eta<1\),
	\[
	u\in C\bigl([0,T];(X_0,X_1)_{\eta,\infty}\bigr).
	\]
	More precisely, for \(0<|t-s|\le T\),
	\[
	\|u(t)-u(s)\|_{(X_0,X_1)_{\eta,\infty}}
	\le
	C_\eta
	\Bigl(|t-s|\log\frac{eT}{|t-s|}\Bigr)^{1-\eta}
	\|u'\|_{BMO_*(0,T;X_0)}^{1-\eta}
	\|u\|_{L^\infty(0,T;X_1)}^\eta.
	\]
	Consequently,
	\[
	u\in C^{0,\alpha}\bigl([0,T];(X_0,X_1)_{\eta,\infty}\bigr)
	\]
	for every \(0<\alpha<1-\eta\).
	\end{theorem}

	\begin{proof}
	We first prove the base-space logarithmic modulus.  Let \(v\in BMO_*(0,T;X_0)\).  We claim that for every interval \(I\subset(0,T)\) of length \(h\),
	\[
	\|v_I\|_{X_0}
	\le
	C\|v\|_{BMO_*(0,T;X_0)}\log\frac{eT}{h}.
	\]
	Indeed, let \(J=(0,T)\).  Enlarge \(I\) dyadically inside \(J\): choose intervals
	\[
	I=I_0\subset I_1\subset\cdots\subset I_N\subset J
	\]
	with \(|I_{k+1}|\le 2|I_k|\) and \(N\le C\log(eT/h)\), and with \(|I_N|\simeq T\).  For nested intervals \(I_k\subset I_{k+1}\),
	\[
	\|v_{I_k}-v_{I_{k+1}}\|_{X_0}
	\le
	\frac1{|I_k|}\int_{I_k}\|v(t)-v_{I_{k+1}}\|_{X_0}\,dt
	\le
	C[v]_{BMO(0,T;X_0)}.
	\]
	Telescoping gives
	\[
	\|v_I\|_{X_0}
	\le
	\|v_J\|_{X_0}+C N[v]_{BMO(0,T;X_0)}
	\le
	C\|v\|_{BMO_*(0,T;X_0)}\log\frac{eT}{h}.
	\]

	Apply this to \(v=u'\).  If \(I=(s,t)\), then
	\[
	u(t)-u(s)=\int_s^t u'(\tau)\,d\tau=|I|(u')_I
	\]
	in \(X_0\).  Hence
	\[
	\|u(t)-u(s)\|_{X_0}
	\le
	C |t-s|\log\frac{eT}{|t-s|}
	\|u'\|_{BMO_*(0,T;X_0)}.
	\]
	On the other hand, for almost every \(s,t\),
	\[
	\|u(t)-u(s)\|_{X_1}
	\le 2\|u\|_{L^\infty(0,T;X_1)}.
	\]
	Using the elementary interpolation inequality
	\[
	\|x\|_{(X_0,X_1)_{\eta,\infty}}
	\le
	C_\eta\|x\|_{X_0}^{1-\eta}\|x\|_{X_1}^{\eta},
	\qquad x\in X_1,
	\]
	with \(x=u(t)-u(s)\), we obtain the stated logarithmic modulus at Lebesgue times.  The representative on all of \([0,T]\) is obtained by completion in \((X_0,X_1)_{\eta,\infty}\), using the modulus estimate and the \(X_0\)-continuous representative.  Finally, since
	\[
	h\log(eT/h)\le C_{\alpha,\eta,T}h^{\alpha/(1-\eta)}
	\]
	for every \(0<h\le T\) and every \(0<\alpha<1-\eta\), the H\"older conclusion follows.
	\end{proof}

	\begin{proposition}\label{prop:bmo-log-sharp}
	The logarithmic factor in Theorem~\ref{thm:bmo-log-trace} cannot in general be removed.
	\end{proposition}

	\begin{proof}
	It is enough to consider the scalar case \(X_0=X_1=\mathbb R\).  On \((0,1)\), let
	\[
	v(t)=\log\frac{e}{t}.
	\]
	The function \(v\) belongs to scalar \(BMO(0,1)\).  This follows directly from the scale invariance of the logarithm near the origin: on intervals \((0,r)\), the mean oscillation is independent of \(r\), and on intervals away from the origin the logarithm has bounded mean oscillation by elementary monotonicity estimates.
	Define
	\[
	u(t)=\int_0^t v(s)\,ds.
	\]
	Then
	\[
	u(t)=t\log\frac{e}{t}+t
		\simeq t\log\frac{e}{t}
	\]
	as \(t\downarrow0\).  Hence \(u'\in BMO(0,1)\), but \(u\) does not satisfy a Lipschitz estimate at the origin.  More precisely,
	\[
	\frac{|u(t)-u(0)|}{t}=\log\frac{e}{t}+1\to\infty.
	\]
	Thus a theorem assuming only \(u'\in BMO\) cannot replace the modulus \(t\log(e/t)\) by \(Ct\).  After interpolation, the corresponding factor \((t\log(e/t))^{1-\eta}\) is therefore the natural endpoint modulus.
	\end{proof}


	
\section{Concluding remarks}
The exact trace space of $\mathbb E^{p,\lambda}$ is the weak interpolation space $(X_0,X_1)_{\theta,\infty}$. The sharpness results show that both the smoothness exponent and the weak fine index are intrinsic to Morrey control. The endpoint estimates provide a complementary bounded-mean-oscillation trace principle and identify the corresponding logarithmic modulus of continuity.

\section*{Statements and Declarations}
\noindent\textbf{Competing interests.} The authors declare no competing interests.

\noindent\textbf{Data availability.} No datasets were generated or analysed during this study.

\noindent\textbf{Use of generative AI.} During preparation of the manuscript, the first author used Gemini for language editing and the literature exposition. All mathematical statements and proofs were reviewed and verified by the authors, who take full responsibility for the content of the manuscript.

\begin{thebibliography}{99}
\bibitem{LionsPeetre1964}
J.-L. Lions and J. Peetre,
\emph{Sur une classe d'espaces d'interpolation},
Publ. Math. Inst. Hautes \'Etudes Sci. \textbf{19} (1964), 5--68.
\url{https://doi.org/10.1007/BF02684796}.

\bibitem{Peetre1966}
J. Peetre,
\emph{On the theory of interpolation spaces},
Rev. Uni\'on Mat. Argentina \textbf{23} (1966/67), 49--66.

\bibitem{Weis2001}
L. Weis,
\emph{Operator-valued Fourier multiplier theorems and maximal $L_p$-regularity},
Math. Ann. \textbf{319} (2001), 735--758.
\url{https://doi.org/10.1007/s002080100177}.

\bibitem{KunstmannWeis2004}
P. C. Kunstmann and L. Weis,
\emph{Maximal $L^p$-regularity for parabolic equations, Fourier multiplier theorems and $H^\infty$-functional calculus},
in \emph{Functional Analytic Methods for Evolution Equations}, Lecture Notes in Math. \textbf{1855}, Springer, Berlin, 2004, 65--311.
\url{https://doi.org/10.1007/978-3-540-44653-8_2}.

\bibitem{Schneider2010}
C. Schneider,
\emph{Trace operators in Besov and Triebel--Lizorkin spaces},
Z. Anal. Anwend. \textbf{29} (2010), 275--302.
\url{https://doi.org/10.4171/ZAA/1409}.

\bibitem{JohnsenSickel2018}
J. Johnsen and W. Sickel,
\emph{On the trace problem for Lizorkin--Triebel spaces with mixed norms},
Math. Nachr. \textbf{281} (2008), 669--696.
\url{https://doi.org/10.1002/mana.200610634}.

\bibitem{MeyriesVeraar2014}
M. Meyries and M. C. Veraar,
\emph{Traces and embeddings of anisotropic function spaces},
Math. Ann. \textbf{360} (2014), 571--606.
\url{https://doi.org/10.1007/s00208-014-1042-6}.

\bibitem{AgrestiLindemulderVeraar2023}
A. Agresti, N. Lindemulder and M. C. Veraar,
\emph{On the trace embedding and its applications to evolution equations},
Math. Nachr. \textbf{296} (2023), 1319--1350.
\url{https://doi.org/10.1002/mana.202100192}.

\bibitem{DenkRoodenburg2026}
R. Denk and F. B. Roodenburg,
\emph{Trace theory for parabolic boundary value problems with rough boundary conditions},
J. Evol. Equ. \textbf{26} (2026), article 80.
\url{https://doi.org/10.1007/s00028-026-01226-6}.

\bibitem{BurenkovGhorbanalizadehSawano2019}
V. I. Burenkov, A. Ghorbanalizadeh and Y. Sawano,
\emph{On the equivalence of the $K$-functional and the modulus of continuity on the Morrey spaces},
J. Approx. Theory \textbf{248} (2019), 105295.
\url{https://doi.org/10.1016/j.jat.2019.105295}.

\bibitem{BurenkovEtAl2010}
V. I. Burenkov and E. D. Nursultanov,
\emph{Description of interpolation spaces for local Morrey-type spaces},
Proc. Steklov Inst. Math. \textbf{269} (2010), 46--56.
\url{https://doi.org/10.1134/S0081543810020045}.

\bibitem{Peetre1969Morrey}
J. Peetre,
\emph{On the theory of $L^{p,\lambda}$ spaces},
J. Funct. Anal. \textbf{4} (1969), 71--87.
\url{https://doi.org/10.1016/0022-1236(69)90022-6}.

\bibitem{JohnNirenberg1961}
F. John and L. Nirenberg,
\emph{On functions of bounded mean oscillation},
Comm. Pure Appl. Math. \textbf{14} (1961), 415--426.
\url{https://doi.org/10.1002/cpa.3160140317}.

\bibitem{Sawano2018}
Y. Sawano,
\emph{Theory of Besov Spaces},
Developments in Mathematics \textbf{56}, Springer, Singapore, 2018.
\url{https://doi.org/10.1007/978-981-13-0836-9}.

\bibitem{OgawaSuguro2023}
T. Ogawa and T. Suguro,
\emph{Maximal regularity of the heat evolution equation on spatial local spaces and application to a singular limit problem of the Keller--Segel system},
Math. Ann. \textbf{387} (2023), 389--431.
\url{https://doi.org/10.1007/s00208-022-02469-7}.
\end{thebibliography}

\end{document}
